%% paper65_graph_theory_en.tex
%% Graph Theory: Foundations, Coloring, Planarity, and the Hadwiger-Nelson Problem
%% Author: Michael Fuhrmann
%% Project: Specialist Mathematics Project, Build 145
%% Date: 2026-03-12

\documentclass[12pt,a4paper]{amsart}

%% ---------------------------------------------------------------
%% Packages
%% ---------------------------------------------------------------
\usepackage{amsmath, amssymb, amsthm}
\usepackage{mathtools}
\usepackage{geometry}
\usepackage{hyperref}
\usepackage{enumitem}
\usepackage{tikz}
\usepackage{xcolor}
\usepackage{booktabs}
\usepackage{microtype}

\geometry{margin=2.5cm}

\hypersetup{
  colorlinks=true,
  linkcolor=blue!60!black,
  citecolor=green!50!black,
  urlcolor=blue!70!black
}

%% ---------------------------------------------------------------
%% Theorem Environments
%% ---------------------------------------------------------------
\newtheorem{theorem}{Theorem}[section]
\newtheorem{lemma}[theorem]{Lemma}
\newtheorem{corollary}[theorem]{Corollary}
\newtheorem{proposition}[theorem]{Proposition}
\theoremstyle{definition}
\newtheorem{definition}[theorem]{Definition}
\newtheorem{example}[theorem]{Example}
\theoremstyle{remark}
\newtheorem{remark}[theorem]{Remark}
\newtheorem{conjecture}[theorem]{Conjecture}

%% ---------------------------------------------------------------
%% Custom Commands
%% ---------------------------------------------------------------
\newcommand{\Z}{\mathbb{Z}}
\newcommand{\R}{\mathbb{R}}
\newcommand{\N}{\mathbb{N}}
\newcommand{\C}{\mathbb{C}}
\newcommand{\chr}{\chi}
\newcommand{\degv}{\mathrm{deg}}

%% ---------------------------------------------------------------
%% Document
%% ---------------------------------------------------------------

\begin{document}
\title{Graph Theory: Foundations, Coloring, Planarity,\\
       and the Hadwiger-Nelson Problem}

\author{Michael Fuhrmann}

\address{Specialist Mathematics Project, Build 145}

\email{specialist-maths@localhost}

\date{2026-03-12}

\keywords{graph theory, graph coloring, planar graphs, spanning trees,
          spectral graph theory, Hadwiger-Nelson problem, four color theorem,
          chromatic number, Euler formula, Menger's theorem}

\subjclass[2020]{05C15, 05C10, 05C40, 05C50, 05C45, 05D10}

\begin{abstract}
Graph theory, born from Euler's solution of the K\"{o}nigsberg bridge problem
in 1736, has grown into one of the most versatile branches of discrete
mathematics, with applications in computer science, physics, chemistry, and
combinatorics. This paper presents a self-contained development of the central
results of graph theory: the Handshake Lemma, characterizations of trees
(including Cayley's formula $\kappa(K_n) = n^{n-2}$), Menger's theorem and the
Max-Flow Min-Cut theorem, graph coloring theory culminating in the celebrated
Four Color Theorem (Appel-Haken 1976; Robertson et al.\ 1997), planarity and
Kuratowski's characterization via $K_5$ and $K_{3,3}$, Euler's formula
$V - E + F = 2$, the Five Color Theorem, spectral graph theory with the
Cheeger inequality, Eulerian and Hamiltonian cycle theory with the NP-completeness
of the Hamiltonian path problem, and extremal graph theory including the
Tur\'{a}n theorem and the Szemer\'{e}di Regularity Lemma. A central focus is the
\emph{Hadwiger-Nelson problem}: how many colors are needed to color the plane
so that no two points at unit distance share a color? We present the classical
bounds $4 \le \chr(\R^2) \le 7$, the Moser spindle and Golomb graph proving
$\chr(\R^2) \ge 4$, and the landmark 2018 result of de Grey establishing
$\chr(\R^2) \ge 5$. The exact value of $\chr(\R^2)$ remains an open problem.
All proofs are presented in full detail; conjectural statements are clearly
distinguished from established theorems.
\end{abstract}

\maketitle

\tableofcontents

%% =============================================================
\section{Introduction}
%% =============================================================

The origins of graph theory lie in a famous puzzle posed about the city of
K\"{o}nigsberg (now Kaliningrad). The Pregel river divided the city into four
regions connected by seven bridges. The question was whether it was possible
to traverse all seven bridges exactly once and return to the starting point.
In 1736, Leonhard Euler proved that no such walk exists, thereby solving the
first problem in graph theory and founding the discipline itself
\cite{euler1736}.

Euler's insight was to abstract the geography away entirely: replace the four
landmasses by \emph{vertices} and the seven bridges by \emph{edges}. The
question then becomes purely combinatorial. He observed that a closed walk
traversing every edge exactly once (today called an \emph{Eulerian circuit})
exists if and only if every vertex has even degree. Since K\"{o}nigsberg
has four vertices of odd degree, no such circuit exists.

From this humble beginning, graph theory has expanded enormously. Today it
encompasses:
\begin{itemize}[leftmargin=2em]
  \item \textbf{Structural theory}: trees, connectivity, planarity, minors;
  \item \textbf{Coloring theory}: proper colorings, chromatic polynomials,
        the four color theorem;
  \item \textbf{Spectral theory}: eigenvalues of adjacency and Laplacian matrices;
  \item \textbf{Extremal theory}: Tur\'{a}n-type problems, Ramsey theory;
  \item \textbf{Algorithmic aspects}: matching, network flow, NP-completeness;
  \item \textbf{Geometric graph theory}: unit-distance graphs, the
        Hadwiger-Nelson problem.
\end{itemize}

This paper is organized as follows. Section~\ref{sec:basics} introduces the
fundamental definitions and the Handshake Lemma. Section~\ref{sec:trees}
treats trees and spanning trees, including Cayley's formula. Section~\ref{sec:connectivity}
covers connectivity, Menger's theorem, and max-flow min-cut. Section~\ref{sec:coloring}
develops graph coloring theory and the Four Color Theorem. Section~\ref{sec:planar}
studies planar graphs, Euler's formula, and Kuratowski's theorem.
Section~\ref{sec:hadwiger} is dedicated to the Hadwiger-Nelson problem.
Section~\ref{sec:spectral} introduces spectral graph theory. Section~\ref{sec:euler_hamilton}
treats Eulerian and Hamiltonian graphs. Section~\ref{sec:extremal} covers
extremal graph theory. Section~\ref{sec:references} collects the references.

\subsection*{Conventions}

Throughout this paper, $G = (V, E)$ denotes a simple undirected graph with
vertex set $V$ and edge set $E \subseteq \binom{V}{2}$. We write $n = |V|$
and $m = |E|$. A \emph{simple} graph has no loops and no multiple edges.
We denote the complete graph on $n$ vertices by $K_n$, the complete bipartite
graph by $K_{p,q}$, the path on $n$ vertices by $P_n$, and the cycle on $n$
vertices by $C_n$.

%% =============================================================
\section{Graph Basics}
\label{sec:basics}
%% =============================================================

\begin{definition}[Graph]
A \emph{simple undirected graph} is a pair $G = (V, E)$ where $V$ is a finite
nonempty set of \emph{vertices} and $E \subseteq \binom{V}{2}$ is a set of
\emph{edges}, each being a two-element subset of $V$. A \emph{directed graph}
(digraph) $\vec{G} = (V, A)$ has a set $A \subseteq V \times V$ of
\emph{arcs} (directed edges), where $(u, v) \ne (v, u)$ in general.
\end{definition}

\begin{definition}[Degree]
The \emph{degree} $\degv(v)$ of a vertex $v \in V$ is the number of edges
incident to $v$:
\[
  \degv(v) = |\{u \in V : \{u,v\} \in E\}|.
\]
The \emph{minimum degree} is $\delta(G) = \min_{v \in V} \degv(v)$ and the
\emph{maximum degree} is $\Delta(G) = \max_{v \in V} \degv(v)$.
\end{definition}

\begin{definition}[Adjacency Matrix]
The \emph{adjacency matrix} of $G$ is the $n \times n$ matrix $A = (a_{ij})$
where $a_{ij} = 1$ if $\{i,j\} \in E$ and $a_{ij} = 0$ otherwise. For
simple graphs, $A$ is symmetric with zero diagonal.
\end{definition}

\begin{theorem}[Handshake Lemma]
\label{thm:handshake}
For any graph $G = (V, E)$,
\[
  \sum_{v \in V} \degv(v) = 2|E|.
\]
In particular, the number of vertices of odd degree is always even.
\end{theorem}

\begin{proof}
Each edge $\{u, v\} \in E$ contributes exactly $1$ to $\degv(u)$ and exactly
$1$ to $\degv(v)$. Therefore it contributes exactly $2$ to the sum
$\sum_{v \in V} \degv(v)$. Summing over all $|E|$ edges gives
$\sum_{v \in V} \degv(v) = 2|E|$.

Since $2|E|$ is even and $\sum_{v \in V} \degv(v)$ is a sum of integers equal
to an even number, the number of odd summands must be even.
\end{proof}

\begin{definition}[Subgraph, Isomorphism]
A graph $H = (W, F)$ is a \emph{subgraph} of $G = (V, E)$ if $W \subseteq V$
and $F \subseteq E$. It is an \emph{induced subgraph} on $S \subseteq V$ if
$F = E \cap \binom{S}{2}$. Two graphs $G$ and $H$ are \emph{isomorphic} if
there exists a bijection $\phi: V(G) \to V(H)$ such that $\{u,v\} \in E(G)$
iff $\{\phi(u), \phi(v)\} \in E(H)$.
\end{definition}

\begin{example}
The complete graph $K_n$ has $n$ vertices, $\binom{n}{2}$ edges, and every
vertex has degree $n-1$. The Handshake Lemma gives
$n(n-1) = 2\binom{n}{2}$, which checks out.
\end{example}

\begin{definition}[Walk, Path, Cycle]
A \emph{walk} in $G$ is a sequence $v_0, e_1, v_1, e_2, \ldots, e_k, v_k$
where each $e_i = \{v_{i-1}, v_i\} \in E$. A \emph{path} is a walk without
repeated vertices. A \emph{cycle} is a closed walk ($v_0 = v_k$) without
other repeated vertices, of length $k \ge 3$.
\end{definition}

%% =============================================================
\section{Trees and Spanning Trees}
\label{sec:trees}
%% =============================================================

\begin{definition}[Tree, Forest, Spanning Tree]
A \emph{forest} is an acyclic graph. A \emph{tree} is a connected acyclic
graph. A \emph{spanning tree} of a connected graph $G$ is a subgraph
$T \subseteq G$ that is a tree and contains all vertices of $G$.
\end{definition}

\begin{theorem}[Characterization of Trees]
\label{thm:tree-char}
For a graph $G = (V, E)$ with $|V| = n \ge 1$, the following are equivalent:
\begin{enumerate}[label=\upshape(\roman*)]
  \item $G$ is a tree (connected and acyclic);
  \item $G$ is connected and $|E| = n - 1$;
  \item $G$ is acyclic and $|E| = n - 1$;
  \item Any two distinct vertices of $G$ are connected by a unique path;
  \item $G$ is connected, but deleting any edge makes it disconnected.
\end{enumerate}
\end{theorem}

\begin{proof}
We prove (i) $\Rightarrow$ (ii) $\Rightarrow$ (iii) $\Rightarrow$ (i) and
sketch the equivalence with (iv) and (v).

\medskip
\noindent\textbf{(i) $\Rightarrow$ (ii).}
By induction on $n$. For $n = 1$: the only tree on one vertex has $0 = 1-1$
edges. For $n \ge 2$: every tree has a leaf (a vertex of degree 1). Let $v$
be a leaf of $T$. Then $T - v$ is a connected acyclic graph on $n - 1$
vertices, so by induction $|E(T-v)| = n - 2$. Since $T$ has one more edge
than $T - v$ (the edge incident to $v$), we get $|E(T)| = n - 1$.

\medskip
\noindent\textbf{(ii) $\Rightarrow$ (iii).}
Suppose $G$ is connected with $n - 1$ edges. We show $G$ is acyclic.
A connected graph with $k$ vertices has at least $k - 1$ edges (achieved only
by trees). If $G$ contained a cycle $C$, removing an edge of $C$ would still
leave $G$ connected, reducing to a connected graph on $n$ vertices with $n-2$
edges. But a connected graph on $n$ vertices needs at least $n-1$ edges —
contradiction.

\medskip
\noindent\textbf{(iii) $\Rightarrow$ (i).}
Suppose $G$ is acyclic with $n - 1$ edges. A forest on $n$ vertices with $k$
connected components has exactly $n - k$ edges. Since $|E| = n - 1$, we get
$k = 1$, so $G$ is connected, hence a tree.

\medskip
\noindent The equivalences with (iv) and (v) follow by observing that unique
paths characterize acyclic connected graphs, and minimality of connectivity
characterizes trees.
\end{proof}

\begin{theorem}[Cayley's Formula]
\label{thm:cayley}
The number of distinct labeled spanning trees of $K_n$ is $n^{n-2}$.
\end{theorem}

\begin{proof}[Proof (Pr\"{u}fer sequences)]
We establish a bijection between labeled spanning trees on $[n] = \{1, \ldots, n\}$
and sequences $(a_1, \ldots, a_{n-2}) \in [n]^{n-2}$.

\textbf{Encoding.} Given a tree $T$, we iteratively remove the leaf with
smallest label: at each step $i$, let $\ell_i$ be the smallest-label leaf,
set $a_i$ to be the unique neighbor of $\ell_i$, and remove $\ell_i$. After
$n - 2$ steps, two vertices remain; record the sequence $(a_1, \ldots, a_{n-2})$.

\textbf{Decoding.} Given a sequence $(a_1, \ldots, a_{n-2})$, we recover the
tree: at step $i$, the current leaf is the smallest element of $[n]$ not in
$\{a_i, a_{i+1}, \ldots, a_{n-2}\}$ and not yet removed; add an edge from
this leaf to $a_i$.

One verifies that encoding and decoding are inverse bijections between the
set of labeled spanning trees of $K_n$ and $[n]^{n-2}$. Since $|[n]^{n-2}| = n^{n-2}$,
the formula follows.
\end{proof}

\begin{remark}
For $n = 4$: $4^2 = 16$ labeled spanning trees of $K_4$. For $n = 2$: $2^0 = 1$
(the single edge). These match direct enumeration.
\end{remark}

%% =============================================================
\section{Connectivity}
\label{sec:connectivity}
%% =============================================================

\begin{definition}[Connectivity, Cut]
A graph $G$ is \emph{$k$-connected} ($k$-vertex-connected) if $|V| > k$ and
removing any $k-1$ vertices leaves $G$ connected. The \emph{vertex connectivity}
$\kappa(G)$ is the maximum such $k$. An \emph{$s$-$t$ vertex cut} is a set
$S \subset V \setminus \{s,t\}$ whose removal disconnects $s$ from $t$.
\end{definition}

\begin{theorem}[Menger's Theorem]
\label{thm:menger}
Let $G$ be a graph and $s, t \in V$ with $\{s,t\} \notin E$. The maximum
number of pairwise internally vertex-disjoint $s$-$t$ paths equals the minimum
size of an $s$-$t$ vertex cut.
\end{theorem}

\begin{proof}[Proof sketch]
Let $f^*$ denote the maximum number of vertex-disjoint $s$-$t$ paths and
$c^*$ the minimum vertex cut size. Clearly $f^* \le c^*$ (any cut must
``block'' each path). For the other direction, we use a flow argument.

Construct a directed network $N$ from $G$: replace each vertex $v \ne s, t$
by two vertices $v^{\text{in}}$ and $v^{\text{out}}$ connected by an arc of
capacity $1$; replace each edge $\{u,v\}$ by arcs $(u^{\text{out}}, v^{\text{in}})$
and $(v^{\text{out}}, u^{\text{in}})$, each of capacity $\infty$. A maximum
integer $s$-$t$ flow in $N$ corresponds to vertex-disjoint paths in $G$,
and a minimum $s$-$t$ cut in $N$ corresponds to a minimum vertex cut in $G$.
By the Max-Flow Min-Cut Theorem (Theorem~\ref{thm:maxflow}), these values
are equal.
\end{proof}

\begin{theorem}[Max-Flow Min-Cut Theorem, Ford-Fulkerson]
\label{thm:maxflow}
In a directed network $N = (V, A, c)$ with source $s$ and sink $t$ and
nonnegative integer capacities $c: A \to \Z_{\ge 0}$, the value of a maximum
$s$-$t$ flow equals the capacity of a minimum $s$-$t$ cut.
\end{theorem}

\begin{proof}
Let $f$ be any $s$-$t$ flow and $(S, T)$ any $s$-$t$ cut (with $s \in S$,
$t \in T$, $S \cup T = V$, $S \cap T = \emptyset$). By flow conservation,
the flow value $|f| = \sum_{(u,v): u \in S, v \in T} f(u,v) - \sum_{(u,v): u \in T, v \in S} f(u,v)
\le \sum_{(u,v): u \in S, v \in T} c(u,v) = c(S,T)$.
So $\max |f| \le \min c(S,T)$.

For equality, run the Ford-Fulkerson algorithm: repeatedly find an
\emph{augmenting path} from $s$ to $t$ in the residual graph and increase
the flow along it. This terminates for integer capacities. When no augmenting
path exists, define $S$ to be the set of vertices reachable from $s$ in the
residual graph. Then $(S, V \setminus S)$ is a cut of capacity $|f|$.
Hence $\max |f| = \min c(S,T)$.
\end{proof}

%% =============================================================
\section{Graph Coloring}
\label{sec:coloring}
%% =============================================================

\begin{definition}[Proper Coloring, Chromatic Number]
A \emph{proper $k$-coloring} of $G$ is a function $c: V \to \{1, \ldots, k\}$
such that $c(u) \ne c(v)$ whenever $\{u,v\} \in E$. The \emph{chromatic
number} $\chr(G)$ is the smallest $k$ for which a proper $k$-coloring exists.
\end{definition}

\begin{example}
$\chr(K_n) = n$, $\chr(C_{2k}) = 2$, $\chr(C_{2k+1}) = 3$, $\chr(P_n) = 2$
for $n \ge 2$.
\end{example}

\begin{proposition}[Greedy Coloring Upper Bound]
\label{prop:greedy}
For any graph $G$, $\chr(G) \le \Delta(G) + 1$.
\end{proposition}

\begin{proof}
Order the vertices arbitrarily as $v_1, \ldots, v_n$. Assign to each vertex
$v_i$ the smallest color in $\{1, 2, \ldots\}$ not used by any of its
already-colored neighbors. Since $v_i$ has at most $\Delta(G)$ neighbors, at
most $\Delta(G)$ colors are excluded, so a color from $\{1, \ldots, \Delta(G)+1\}$
is always available.
\end{proof}

\begin{theorem}[Brooks' Theorem]
\label{thm:brooks}
If $G$ is a connected graph that is neither a complete graph nor an odd cycle,
then $\chr(G) \le \Delta(G)$.
\end{theorem}

\begin{proof}[Proof sketch]
We induct on $|V|$. For $\Delta(G) = 2$, $G$ must be a path or an even cycle
(odd cycles and triangles are excluded), both 2-colorable.

For $\Delta(G) \ge 3$: if $G$ is not $\Delta$-regular, remove a vertex $v$
of degree $\delta(G) < \Delta(G)$. The remaining graph $G - v$ can be
$\Delta$-colored by induction, and since $v$ has fewer than $\Delta$ neighbors,
a color remains available for $v$.

If $G$ is $\Delta$-regular with a cut vertex, decompose and apply induction.
If $G$ is $\Delta$-regular and 2-connected, a careful reordering of vertices
(using a path from a vertex and back) allows the greedy algorithm to succeed
with $\Delta$ colors.
\end{proof}

\begin{theorem}[Four Color Theorem]
\label{thm:4color}
Every planar graph is $4$-colorable: $\chr(G) \le 4$ for all planar graphs $G$.
\end{theorem}

\begin{proof}[Historical note and proof outline]
The Four Color Theorem has a rich history. Francis Guthrie conjectured it in
1852. Alfred Kempe published a claimed proof in 1879, which Percy Heawood
refuted in 1890 (finding a gap), though Kempe's argument still yielded the
Five Color Theorem (Theorem~\ref{thm:5color}).

The theorem was finally \textbf{proved by Appel and Haken in 1976}
\cite{appelhaken1977}, using an unavoidable set of 1936 reducible
configurations verified by computer. This was the first major theorem whose
proof relied essentially on computer assistance.

A second, cleaner proof was given by \textbf{Robertson, Sanders, Seymour, and
Thomas in 1997} \cite{robertson1997}, reducing the unavoidable set to 633
configurations and making the computer verification more transparent.

\textbf{Proof strategy.} The argument proceeds in two steps:
\begin{enumerate}
  \item \textbf{Unavoidability:} Show that every planar graph contains at least
        one configuration from a finite list (e.g., configurations of patches
        with small faces around a vertex).
  \item \textbf{Reducibility:} Show that if any such configuration appears
        inside a minimal counterexample, then a smaller planar graph not
        4-colorable can be derived — contradicting minimality.
\end{enumerate}
The combination of unavoidability and reducibility shows no minimal
counterexample can exist, hence no counterexample at all.
\end{proof}

\begin{remark}
No elementary (proof-by-hand) proof of the Four Color Theorem is known. Both
known proofs are computer-assisted.
\end{remark}

%% =============================================================
\section{Planar Graphs}
\label{sec:planar}
%% =============================================================

\begin{definition}[Planar Graph, Embedding, Face]
A graph $G$ is \emph{planar} if it can be drawn in the plane $\R^2$ with no
two edges crossing except at shared endpoints. Such a drawing is a \emph{planar
embedding}. A planar embedding divides the plane into connected regions called
\emph{faces}, including exactly one unbounded face.
\end{definition}

\begin{theorem}[Euler's Formula]
\label{thm:euler-formula}
Let $G$ be a connected planar graph with $V$ vertices, $E$ edges, and $F$ faces
in a planar embedding. Then
\[
  V - E + F = 2.
\]
\end{theorem}

\begin{proof}
By induction on the number of edges. If $G$ is a tree, then $E = V - 1$,
$F = 1$ (only the unbounded face), and $V - (V-1) + 1 = 2$. \checkmark

Now suppose $G$ has a cycle, so there exists an edge $e$ on a cycle. Removing
$e$ merges two faces into one, so $F' = F - 1$, $E' = E - 1$, $V' = V$.
By induction, $V' - E' + F' = 2$, i.e., $V - (E-1) + (F-1) = 2$, which
gives $V - E + F = 2$. \checkmark

If $G$ is connected but has a bridge $e$ (not on any cycle), removing $e$
disconnects $G$: $V' = V$, $E' = E - 1$. The two resulting components
$G_1, G_2$ satisfy $V_1 - E_1 + F_1 = 2$ and $V_2 - E_2 + F_2 = 2$. The
original embedding counted the shared face once, so $F = F_1 + F_2 - 1$
and $V = V_1 + V_2$, $E = E_1 + E_2 + 1$. Adding the two equations and
adjusting gives $V - E + F = 2$.
\end{proof}

\begin{corollary}[Edge Bound for Planar Graphs]
\label{cor:planar-edges}
If $G$ is a simple planar graph with $n \ge 3$ vertices and $m$ edges, then
$m \le 3n - 6$. If additionally $G$ is triangle-free, then $m \le 2n - 4$.
\end{corollary}

\begin{proof}
Each face is bounded by at least 3 edges. Since each edge borders exactly 2
faces, $3F \le 2E$, so $F \le 2E/3$. Substituting into Euler's formula:
$2 = V - E + F \le n - E + 2E/3 = n - E/3$, giving $E \le 3n - 6$.
For triangle-free graphs, each face has at least 4 edges on its boundary,
giving $4F \le 2E$, hence $E \le 2n - 4$.
\end{proof}

\begin{corollary}
$K_5$ and $K_{3,3}$ are not planar.
\end{corollary}

\begin{proof}
For $K_5$: $n = 5$, $m = 10 > 3 \cdot 5 - 6 = 9$. Violates
Corollary~\ref{cor:planar-edges}.
For $K_{3,3}$: $n = 6$, $m = 9$. It is bipartite, hence triangle-free. But
$9 > 2 \cdot 6 - 4 = 8$. Violates the triangle-free bound.
\end{proof}

\begin{theorem}[Kuratowski's Theorem]
\label{thm:kuratowski}
A graph $G$ is planar if and only if it contains no subdivision of $K_5$ or
$K_{3,3}$ as a subgraph.
\end{theorem}

\begin{proof}[Proof sketch]
The necessity (planar $\Rightarrow$ no $K_5$ or $K_{3,3}$ subdivision) follows
from Corollary~\ref{cor:planar-edges} applied to subdivisions.

For sufficiency, one proceeds by contradiction: assume $G$ is a minimal
non-planar graph. Using properties of minimal non-planar graphs (they are
3-connected; every edge is in a cycle), one shows that $G$ must contain a
$K_5$ or $K_{3,3}$ subdivision. The full proof appears in \cite{diestel2017}.
\end{proof}

\begin{theorem}[Five Color Theorem]
\label{thm:5color}
Every planar graph is $5$-colorable.
\end{theorem}

\begin{proof}
By induction on $n = |V|$. The base cases are trivial. For the inductive step,
let $G$ be a planar graph with $n$ vertices. By Corollary~\ref{cor:planar-edges},
$G$ has a vertex $v$ with $\degv(v) \le 5$.

\textbf{Case 1: $\degv(v) \le 4$.} Remove $v$, apply induction to get a
5-coloring of $G - v$. Since $v$ has at most 4 neighbors, at least one of the
5 colors is free for $v$.

\textbf{Case 2: $\degv(v) = 5$.} Let the neighbors be $v_1, \ldots, v_5$ in
cyclic order around $v$ in the embedding. If two non-adjacent neighbors share
a color in the $(G-v)$-coloring, recolor $v$ with the freed color.

Otherwise, $v_1, \ldots, v_5$ use all 5 distinct colors. Consider the
$v_1$-$v_3$ subgraph induced by colors 1 and 3. If $v_1$ and $v_3$ lie in
different connected components of this 2-colored subgraph, swap the colors in
$v_1$'s component. Now $v_1$ has color 3, freeing color 1 for $v$.

If $v_1$ and $v_3$ are in the same component, there is a path $P$ from $v_1$
to $v_3$ using only colors 1 and 3. Together with the edges $v v_1$ and $v v_3$,
this path $P$ separates $v_2$ from $v_4, v_5$. Then $v_2$ and $v_4$ lie in
different components of the 2-colored subgraph using colors 2 and 4, so we
can swap to free color 2 for $v$.
\end{proof}

%% =============================================================
\section{The Hadwiger-Nelson Problem}
\label{sec:hadwiger}
%% =============================================================

\begin{definition}[Unit-Distance Graph, Chromatic Number of the Plane]
The \emph{unit-distance graph of the plane} $G_{\R^2, 1}$ has vertex set
$\R^2$ and an edge between any two points at Euclidean distance $1$. The
\emph{chromatic number of the plane} is
\[
  \chr(\R^2) = \chr(G_{\R^2, 1})
  = \min\bigl\{k : \exists\; k\text{-coloring of } \R^2 \text{ with no two
    unit-distance points sharing a color}\bigr\}.
\]
\end{definition}

The \emph{Hadwiger-Nelson problem}, posed independently by Hugo Hadwiger,
Edward Nelson, and John Isbell around 1950, asks for the exact value of
$\chr(\R^2)$.

\begin{theorem}[Classical Bounds]
\label{thm:hn-classical}
$4 \le \chr(\R^2) \le 7.$
\end{theorem}

\begin{proof}
\textbf{Upper bound $\chr(\R^2) \le 7$.}
Tile the plane with regular hexagons of side length $s = 0.48 < 1/\sqrt{3}
\approx 0.577$. Color the hexagons with 7 colors using the periodic pattern
$(i,j) \mapsto (i + 2j) \bmod 7$ in the hexagonal lattice coordinates.

Two points in the same hexagon have distance at most $2s = 0.96 < 1$. Two
points in same-colored hexagons have distance at least $\sqrt{7} \cdot s\sqrt{3}
= \sqrt{21} \cdot 0.48 \approx 2.20 > 1$. Hence same-colored points always
have distance $\ne 1$, proving $\chr(\R^2) \le 7$.

\textbf{Lower bound $\chr(\R^2) \ge 4$.}
The \emph{Moser spindle} (Moser \& Moser, 1961) is a unit-distance graph
with 7 vertices and 11 edges (see Figure below). Let the vertices be
$A, B, C, D, E, F, G$ embedded in $\R^2$ with all edges having length 1.
The graph has the property that $A$ and $G$ are connected by a path
$A - B - C - D$ and another path $A - E - F - G$, with $D$ and $G$
sharing an edge. A case analysis shows that any 3-coloring forces $A$ and
$G$ to have the same color while $\{A, G\} \in E$ — contradiction. Hence
$\chr(\text{Moser spindle}) = 4 \le \chr(\R^2)$.
\end{proof}

\subsection{The Moser Spindle}

The Moser spindle is constructed as follows. Let:
\begin{align*}
  A &= (0, 0), \quad B = (1, 0), \quad C = (1/2,\, \sqrt{3}/2),\\
  D &= (3/2,\, \sqrt{3}/2), \quad E = (\cos\theta,\, \sin\theta)
      \text{ with } \cos\theta = 5/6,\\
  F &= \bigl((5-\sqrt{33})/12,\, (\sqrt{11}+5\sqrt{3})/12\bigr),\\
  G &= \bigl((15-\sqrt{33})/12,\, (3\sqrt{11}+5\sqrt{3})/12\bigr).
\end{align*}
The edges are: $AB, AC, BC$ (triangle $ABC$), $BD, CD$ ($D$ adjacent to $B$
and $C$), $AE, AF, EF$ (triangle $AEF$), $EG, FG$ ($G$ adjacent to $E$ and
$F$), and $DG$.

One verifies numerically that all 11 edges have length exactly 1.

\begin{lemma}[Chromatic Number of the Moser Spindle]
$\chr(\text{Moser spindle}) = 4$.
\end{lemma}

\begin{proof}
The graph is not 3-colorable. Suppose for contradiction that a valid
3-coloring with colors $\{1, 2, 3\}$ exists. Since $ABC$ is a triangle,
$A$, $B$, $C$ receive three distinct colors; say $c(A) = 1$, $c(B) = 2$,
$c(C) = 3$. Since $BD$ and $CD$ are edges, $c(D) \ne 2$ and $c(D) \ne 3$,
so $c(D) = 1 = c(A)$. Since $AEF$ is a triangle, $c(E), c(F) \in \{2, 3\}$
with $c(E) \ne c(F)$. Since $EG$ and $FG$ are edges, $c(G) \ne c(E)$ and
$c(G) \ne c(F)$, so $c(G)$ must differ from both colors in $\{2,3\}$, forcing
$c(G) = 1$. But $DG$ is an edge and $c(D) = c(G) = 1$: contradiction.
Therefore $\chr(\text{Moser spindle}) \ge 4$. The graph is 4-colorable (e.g.,
assign colors $1,2,3,1,2,3,4$ to $A,\ldots,G$ respecting the edges), so
$\chr = 4$.
\end{proof}

\subsection{De Grey's 2018 Result}

\begin{theorem}[De Grey, 2018]
\label{thm:degrey}
$\chr(\R^2) \ge 5$.
\end{theorem}

\begin{proof}[Proof description]
Aubrey de Grey constructed a finite unit-distance graph $H$ with
$\chr(H) = 5$ \cite{degrey2018}. The construction began with the Moser spindle
and applied a sequence of operations (rotation, union, taking subgraphs) to
obtain graphs that require more and more colors. After substantial computer
search, de Grey found a graph with $\mathbf{1581}$ vertices and $\mathbf{12}$,\!$\mathbf{066}$
edges that is not 4-colorable. This was subsequently verified by multiple
independent teams, and the result was peer-reviewed and published in
\emph{Geombinatorics} in 2019.

Since $H$ is a unit-distance graph in $\R^2$, any coloring of $\R^2$ induces
a coloring of $H$. If $\chr(\R^2) \le 4$, then $H$ would be 4-colorable —
contradiction. Hence $\chr(\R^2) \ge 5$.
\end{proof}

\begin{conjecture}[Exact Value of $\chr(\R^2)$]
\label{conj:hn-exact}
The chromatic number of the plane satisfies $\chr(\R^2) \in \{5, 6, 7\}$.
The exact value is unknown. De Grey's result rules out $\chr(\R^2) \le 4$,
and the hexagonal coloring scheme shows $\chr(\R^2) \le 7$, but whether
$\chr(\R^2) = 5$, $6$, or $7$ remains an open problem.
\end{conjecture}

\begin{remark}
The problem is remarkable because the answer does not depend on the axiom of
choice: any coloring witnessing $\chr(\R^2) = k$ can be taken to be measurable,
and indeed the problem makes sense in a purely combinatorial setting.
The Polymath project (2018) subsequently reduced the number of vertices in a
5-chromatic unit-distance graph to 553 vertices.
\end{remark}

%% =============================================================
\section{Spectral Graph Theory}
\label{sec:spectral}
%% =============================================================

\begin{definition}[Laplacian Matrix]
The \emph{degree matrix} of $G$ is $D = \mathrm{diag}(\degv(v_1), \ldots, \degv(v_n))$.
The \emph{Laplacian matrix} is
\[
  L = D - A,
\]
where $A$ is the adjacency matrix. Explicitly, $L_{ij} = -1$ if $\{i,j\} \in E$,
$L_{ii} = \degv(v_i)$, and $L_{ij} = 0$ otherwise.
\end{definition}

\begin{proposition}[Properties of the Laplacian]
\label{prop:laplacian}
Let $G$ be a graph with $n$ vertices and Laplacian $L$. Then:
\begin{enumerate}[label=\upshape(\roman*)]
  \item $L$ is symmetric and positive semi-definite;
  \item $L \mathbf{1} = \mathbf{0}$, so $0$ is always an eigenvalue with eigenvector
        $\mathbf{1} = (1,\ldots,1)^\top$;
  \item The number of zero eigenvalues of $L$ equals the number of connected
        components of $G$;
  \item For any vector $x \in \R^n$:
        $x^\top L x = \sum_{\{i,j\} \in E} (x_i - x_j)^2 \ge 0$.
\end{enumerate}
\end{proposition}

\begin{proof}
(iv) is a direct computation:
$x^\top L x = x^\top D x - x^\top A x = \sum_i \degv(v_i) x_i^2
- 2\sum_{\{i,j\} \in E} x_i x_j = \sum_{\{i,j\} \in E}(x_i^2 - 2x_ix_j + x_j^2)
= \sum_{\{i,j\} \in E}(x_i - x_j)^2 \ge 0$.
This proves (i). (ii) follows from $\sum_j L_{ij} = \degv(v_i) - \sum_{j: \{i,j\} \in E} 1 = 0$.
(iii) follows because the kernel of $L$ corresponds to locally constant functions
on $G$, i.e., functions constant on each connected component.
\end{proof}

\begin{definition}[Algebraic Connectivity, Fiedler Value]
The eigenvalues of $L$ are $0 = \lambda_1 \le \lambda_2 \le \cdots \le \lambda_n$.
For a connected graph, $\lambda_2 > 0$ and is called the \emph{algebraic
connectivity} or \emph{Fiedler value} of $G$, denoted $a(G)$. The corresponding
eigenvector is the \emph{Fiedler vector}.
\end{definition}

\begin{definition}[Edge Expansion, Cheeger Constant]
The \emph{edge expansion} or \emph{Cheeger constant} of $G$ is
\[
  h(G) = \min_{\emptyset \ne S \subsetneq V,\, |S| \le n/2}
         \frac{|E(S, V \setminus S)|}{|S|},
\]
where $E(S, V\setminus S)$ is the set of edges with one endpoint in $S$ and
the other in $V \setminus S$.
\end{definition}

\begin{theorem}[Cheeger Inequality]
\label{thm:cheeger}
For any $d$-regular graph $G$,
\[
  \frac{h(G)^2}{2d} \le d - \lambda_{n-1}(A) \approx \lambda_2(L)
  \le 2h(G),
\]
where $\lambda_{n-1}(A)$ is the second-largest eigenvalue of the adjacency
matrix $A$. (Equivalently, using the Laplacian: $h(G)^2/(2d) \le \lambda_2(L)/2 \le h(G)$.)
\end{theorem}

\begin{proof}[Proof sketch]
The upper bound $\lambda_2(L) \le 2h(G)$ follows from the Rayleigh quotient
characterization of $\lambda_2$: for any $x$ orthogonal to $\mathbf{1}$,
$\lambda_2 \le x^\top L x / x^\top x$. Choose $x$ to be the characteristic
function of the optimal set $S$ centered at $0$; the resulting bound gives
$\lambda_2 \le 2h(G)$.

The lower bound $h(G) \le \sqrt{2d \lambda_2(L)}$ uses the co-area formula
(discrete Federer-Fleming theorem) to relate the edge boundary of level sets
of the Fiedler vector to the algebraic connectivity.
\end{proof}

\begin{remark}
The Cheeger inequality is the discrete analogue of the isoperimetric inequality
in Riemannian geometry. It connects combinatorial expansion (useful for rapid
mixing of random walks) to the spectral gap.
\end{remark}

%% =============================================================
\section{Eulerian and Hamiltonian Graphs}
\label{sec:euler_hamilton}
%% =============================================================

\begin{definition}[Eulerian Circuit, Trail, Hamiltonian Cycle]
An \emph{Eulerian trail} is a walk that traverses every edge exactly once.
An \emph{Eulerian circuit} is a closed Eulerian trail. A graph containing an
Eulerian circuit is called \emph{Eulerian}. A \emph{Hamiltonian cycle} is a
cycle that visits every vertex exactly once. A graph containing a Hamiltonian
cycle is called \emph{Hamiltonian}.
\end{definition}

\begin{theorem}[Euler's Characterization of Eulerian Graphs]
\label{thm:eulerian}
A connected graph $G$ is Eulerian if and only if every vertex has even degree.
\end{theorem}

\begin{proof}
\textbf{Necessity.} In any Eulerian circuit, each time the trail visits a
vertex $v$, it enters via one edge and exits via another, contributing $2$ to
$\degv(v)$. Hence $\degv(v)$ is even for all $v$.

\textbf{Sufficiency.} Let $G$ be connected with all degrees even. Start a
trail at any vertex $v_0$. The trail can always continue (since even degrees
ensure that arrival at a vertex implies the existence of an unused departure
edge) until it closes at $v_0$. If this closed trail does not use all edges,
there exists an unused edge $e = \{u, w\}$ incident to a vertex $u$ on the
trail (since $G$ is connected). Start a new closed trail from $u$ using only
unused edges; splice this into the original trail at $u$. Repeat until all
edges are used.
\end{proof}

\begin{remark}
For the K\"{o}nigsberg bridge problem: the four landmasses have odd degrees
(3, 3, 3, 5), so no Eulerian circuit exists. This confirms Euler's 1736 result.
\end{remark}

\begin{theorem}[Dirac's Theorem]
\label{thm:dirac}
If $G$ is a simple graph on $n \ge 3$ vertices with $\delta(G) \ge n/2$,
then $G$ is Hamiltonian.
\end{theorem}

\begin{proof}
Let $P = v_1, v_2, \ldots, v_k$ be a longest path in $G$. We claim $k = n$
and the path can be closed into a Hamiltonian cycle.

If $k < n$, then $v_1$ has $\deg(v_1) \ge n/2$ neighbors, all on $P$ (since
$P$ is maximal). Similarly for $v_k$. By a counting argument, there exist
indices $i$ such that $v_k v_i \in E$ and $v_1 v_{i+1} \in E$. This allows
us to form a cycle $C = v_1 \cdots v_i v_k v_{k-1} \cdots v_{i+1} v_1$.
Since $G$ is connected and $k < n$, some vertex $w \notin V(C)$ is adjacent
to a vertex in $C$, allowing extension of $P$ — contradicting maximality.
Hence $k = n$ and $G$ is Hamiltonian.
\end{proof}

\begin{theorem}[NP-Completeness of the Hamiltonian Path Problem]
\label{thm:ham-np}
The problem of deciding whether a given graph $G$ contains a Hamiltonian path
(or cycle) is NP-complete.
\end{theorem}

\begin{proof}
The problem is in NP: a Hamiltonian path, if given, can be verified in
polynomial time (check that each vertex appears exactly once and all edges
are in $E$).

NP-hardness is shown by reduction from 3-SAT. Given a 3-CNF formula
$\phi$ on $n$ variables and $m$ clauses, construct a graph $G_\phi$ such
that $G_\phi$ has a Hamiltonian cycle iff $\phi$ is satisfiable. The standard
reduction (Karp, 1972 \cite{karp1972}) uses a ``variable gadget'' for each
variable (a pair of paths representing true/false) and a ``clause gadget''
for each clause (a single vertex connected to the variable paths). The
resulting graph has $O(nm)$ vertices and edges, constructible in polynomial
time. The correctness of the reduction ensures the equivalence.
\end{proof}

\begin{remark}
The \emph{Travelling Salesman Problem} (TSP) asks for the minimum-weight
Hamiltonian cycle in a weighted complete graph. TSP is NP-hard and a central
problem in combinatorial optimization. For Euclidean TSP (points in the plane,
Euclidean distances), a $(1+\varepsilon)$-approximation exists in polynomial
time for any $\varepsilon > 0$ (PTAS, Arora 1998), but exact solution remains
NP-hard.
\end{remark}

%% =============================================================
\section{Extremal Graph Theory}
\label{sec:extremal}
%% =============================================================

Extremal graph theory asks: how many edges can a graph on $n$ vertices have
while avoiding a certain subgraph?

\begin{definition}[Tur\'{a}n Graph]
The \emph{Tur\'{a}n graph} $T(n, r)$ is the complete $r$-partite graph on
$n$ vertices whose parts are as equal as possible (sizes $\lfloor n/r \rfloor$
or $\lceil n/r \rceil$). Its number of edges is
\[
  t(n, r) = \left(1 - \frac{1}{r}\right)\frac{n^2}{2}
            + O(n).
\]
\end{definition}

\begin{theorem}[Tur\'{a}n's Theorem]
\label{thm:turan}
If $G$ is a simple graph on $n$ vertices with more than $t(n, r)$ edges,
then $G$ contains $K_{r+1}$ as a subgraph. Conversely, $T(n,r)$ is the
unique densest $K_{r+1}$-free graph on $n$ vertices.
\end{theorem}

\begin{proof}
We present the elegant \textbf{Zykov symmetrization} proof.

Call a graph \emph{extremal} if it is $K_{r+1}$-free and has the maximum
number of edges. We claim every extremal graph is $T(n,r)$.

\textbf{Step 1: Extremal graphs are complete multipartite.}
Suppose $G$ is extremal. If two non-adjacent vertices $u, v$ have
$\deg(u) < \deg(v)$, replace $u$ by a copy of $v$ (same neighborhood as $v$).
This does not create $K_{r+1}$ (since $v$'s neighborhood is $K_{r+1}$-free),
and it cannot decrease the edge count. Hence we may assume: for all non-adjacent
$u,v$, $\deg(u) = \deg(v)$.

Non-adjacency is an equivalence relation (if $u \not\sim v$ and $u \not\sim w$,
they must have equal degree, so $v \not\sim w$ as well by maximality). Hence
$G$ is complete multipartite.

\textbf{Step 2: Parts must be balanced.}
If $G = K_{n_1, \ldots, n_r}$ with $n_i > n_j + 1$ for some $i \ne j$, move
a vertex from part $i$ to part $j$. This increases the number of edges (the
change in edges is $n_j - n_i + 1 > 0$ if computed carefully), contradicting
extremality. Hence all parts have size $\lfloor n/r\rfloor$ or $\lceil n/r\rceil$,
i.e., $G = T(n,r)$.
\end{proof}

\begin{theorem}[Zarankiewicz Problem, K\H{o}v\'{a}ri-S\'{o}s-Tur\'{a}n]
\label{thm:zarankiewicz}
If $G$ is a bipartite graph on parts of size $m$ and $n$ (with $m \le n$)
containing no $K_{s,t}$ as a subgraph (with $s \le t$), then
\[
  |E(G)| \le \tfrac{1}{2}\bigl((t-1)^{1/s} n m^{1-1/s} + (s-1)m\bigr).
\]
In particular, ex$(n, K_{2,2}) = O(n^{3/2})$.
\end{theorem}

\begin{proof}[Proof sketch]
Count pairs $(u, \{v_1, v_2, \ldots, v_s\})$ where $u$ is in the $m$-side
and the $v_i$ are in the $n$-side with all $v_i$ adjacent to $u$. On the
one hand, this count is $\sum_{u} \binom{\deg(u)}{s} \ge m \cdot \binom{|E|/m}{s}$
by convexity. On the other hand, for each $s$-set $\{v_1, \ldots, v_s\}$ in
the $n$-side, there are at most $t-1$ common neighbors in the $m$-side
(else $K_{s,t}$ exists). So the count is at most $(t-1)\binom{n}{s}$.
Combining and solving for $|E|$ gives the result.
\end{proof}

\begin{theorem}[Szemer\'{e}di Regularity Lemma]
\label{thm:szemeredi}
For every $\varepsilon > 0$ and integer $k_0 \ge 1$, there exists $M = M(\varepsilon, k_0)$
such that the vertex set of every graph $G$ with $n \ge M$ vertices can be
partitioned into $k$ classes $V_1, \ldots, V_k$ with $k_0 \le k \le M$ and
$|V_1| \le |V_2| \le \cdots \le |V_k| \le |V_1| + 1$, such that all but at
most $\varepsilon k^2$ of the pairs $(V_i, V_j)$ are $\varepsilon$-regular.
\end{theorem}

\begin{proof}[Proof sketch]
The \textbf{energy increment argument}: define the \emph{mean square density}
(energy) $q(P) = \sum_{i<j} \frac{|V_i||V_j|}{n^2} d(V_i, V_j)^2$ for a
partition $P$, where $d(V_i, V_j) = e(V_i, V_j)/(|V_i||V_j|)$ is the edge
density. Note $0 \le q(P) \le 1$.

If the current partition has too many irregular pairs, refine it by splitting
each class using the ``witness'' of irregularity. Each refinement step
increases the energy by at least $\varepsilon^5$. Since energy is bounded by
1, after at most $\varepsilon^{-5}$ refinements the process terminates with
a regular partition.
\end{proof}

\begin{remark}
The Szemer\'{e}di Regularity Lemma is a cornerstone of modern combinatorics.
It enables the \emph{counting lemma} and \emph{removal lemma} and was
central to the proof of Szemer\'{e}di's theorem on arithmetic progressions.
The bound $M(\varepsilon)$ grows as an exponential tower of height $O(1/\varepsilon^5)$,
which is known to be essentially tight (Gowers, 1997).
\end{remark}

%% =============================================================
\section{References}
\label{sec:references}
%% =============================================================

\begin{thebibliography}{99}

\bibitem{euler1736}
L.~Euler,
\emph{Solutio problematis ad geometriam situs pertinentis},
Commentarii Academiae Scientiarum Imperialis Petropolitanae \textbf{8} (1736),
128--140.

\bibitem{cayley1889}
A.~Cayley,
\emph{A theorem on trees},
Quart.\ J.\ Math.\ \textbf{23} (1889), 376--378.

\bibitem{appelhaken1977}
K.~Appel and W.~Haken,
\emph{Every planar map is four colorable. Part~I: Discharging},
Illinois J.\ Math.\ \textbf{21} (1977), 429--490.
K.~Appel, W.~Haken, and J.~Koch,
\emph{Every planar map is four colorable. Part~II: Reducibility},
Illinois J.\ Math.\ \textbf{21} (1977), 491--567.

\bibitem{robertson1997}
N.~Robertson, D.~P.~Sanders, P.~D.~Seymour, and R.~Thomas,
\emph{The four-colour theorem},
J.\ Combin.\ Theory Ser.\ B \textbf{70} (1997), 2--44.

\bibitem{degrey2018}
A.~D.~N.~J.\ de~Grey,
\emph{The chromatic number of the plane is at least 5},
Geombinatorics \textbf{28} (2018/2019), 18--31.
arXiv:1804.02385.

\bibitem{diestel2017}
R.~Diestel,
\emph{Graph Theory},
5th ed., Springer, Berlin, 2017.

\bibitem{bollobas1998}
B.~Bollob\'{a}s,
\emph{Modern Graph Theory},
Springer, New York, 1998.

\bibitem{karp1972}
R.~M.~Karp,
\emph{Reducibility among combinatorial problems},
in: R.~E.~Miller and J.~W.~Thatcher (eds.),
\emph{Complexity of Computer Computations},
Plenum Press, New York, 1972, pp.~85--103.

\bibitem{turan1941}
P.~Tur\'{a}n,
\emph{On an extremal problem in graph theory} (in Hungarian),
Mat.\ Fiz.\ Lapok \textbf{48} (1941), 436--452.

\bibitem{szemeredi1978}
E.~Szemer\'{e}di,
\emph{Regular partitions of graphs},
Probl\`{e}mes Combinatoires et Th\'{e}orie des Graphes,
Colloques Internationaux du CNRS \textbf{260}, Paris, 1978, pp.~399--401.

\bibitem{mosermoser1961}
L.~Moser and W.~Moser,
\emph{Solution to problem 10},
Canad.\ Math.\ Bull.\ \textbf{4} (1961), 187--189.

\bibitem{kuratowski1930}
K.~Kuratowski,
\emph{Sur le probl\`{e}me des courbes gauches en topologie},
Fund.\ Math.\ \textbf{15} (1930), 271--283.

\end{thebibliography}

\end{document}
